\documentclass[11pt]{article}
\newcommand{\ListaTitulo}{Lista 07 --- Contrastes e Comparações Múltiplas}
% Preâmbulo compartilhado — MATD48, Listas de Exercícios 2026
% Usado por Lista{NN}.tex e Gabarito{NN}.tex via % Preâmbulo compartilhado — MATD48, Listas de Exercícios 2026
% Usado por Lista{NN}.tex e Gabarito{NN}.tex via % Preâmbulo compartilhado — MATD48, Listas de Exercícios 2026
% Usado por Lista{NN}.tex e Gabarito{NN}.tex via \input{preamble}

\usepackage[utf8]{inputenc}
\usepackage[T1]{fontenc}
\usepackage[brazil]{babel}
\usepackage{fullpage}
\usepackage{amsmath,amssymb}
\usepackage{enumitem}
\usepackage{xcolor}
\usepackage{fancyhdr}
\usepackage{titlesec}
\usepackage{listings}
\usepackage{hyperref}
\usepackage{booktabs}
\usepackage{graphicx}

\definecolor{matd48azul}{HTML}{0B4F6C}
\definecolor{matd48verde}{HTML}{2E8B57}

\titleformat{\section}{\normalfont\Large\bfseries\color{matd48azul}}{\thesection}{1em}{}
\titleformat{\subsection}{\normalfont\large\bfseries\color{matd48azul}}{\thesubsection}{1em}{}

\providecommand{\ListaTitulo}{Lista de Exercícios}

\setlength{\headheight}{14pt}
\pagestyle{fancy}
\fancyhf{}
\lhead{\small MATD48 --- Planejamento de Experimentos A}
\rhead{\small \ListaTitulo}
\cfoot{\thepage}
\renewcommand{\headrulewidth}{0.4pt}

\lstset{
  basicstyle=\ttfamily\small,
  breaklines=true,
  frame=single,
  columns=fullflexible,
  backgroundcolor=\color{gray!8},
  commentstyle=\color{matd48verde},
  keywordstyle=\color{matd48azul}\bfseries,
  language=R,
  extendedchars=true,
  literate=%
    {á}{{\'a}}1 {à}{{\`a}}1 {â}{{\^a}}1 {ã}{{\~a}}1
    {é}{{\'e}}1 {ê}{{\^e}}1 {í}{{\'i}}1 {ó}{{\'o}}1
    {ô}{{\^o}}1 {õ}{{\~o}}1 {ú}{{\'u}}1 {ç}{{\c c}}1
    {Á}{{\'A}}1 {À}{{\`A}}1 {Â}{{\^A}}1 {Ã}{{\~A}}1
    {É}{{\'E}}1 {Ê}{{\^E}}1 {Í}{{\'I}}1 {Ó}{{\'O}}1
    {Ô}{{\^O}}1 {Õ}{{\~O}}1 {Ú}{{\'U}}1 {Ç}{{\c C}}1
}

\hypersetup{colorlinks=true, linkcolor=matd48azul, urlcolor=matd48azul, citecolor=matd48azul}

\newcommand{\cabecalho}[2]{%
  \begin{center}
    {\Large\bfseries\color{matd48azul} #1}\\[0.3em]
    {\large #2}\\[0.3em]
    \rule{\linewidth}{0.6pt}
  \end{center}
  \vspace{0.5em}
}

\newenvironment{questao}[1]{\vspace{1em}\noindent\textbf{Questão #1.}\hspace{0.4em}}{\par}
\newenvironment{solucao}{\vspace{0.3em}\noindent\textit{Solução.}\hspace{0.4em}\color{black!85}}{\par\vspace{0.5em}}


\usepackage[utf8]{inputenc}
\usepackage[T1]{fontenc}
\usepackage[brazil]{babel}
\usepackage{fullpage}
\usepackage{amsmath,amssymb}
\usepackage{enumitem}
\usepackage{xcolor}
\usepackage{fancyhdr}
\usepackage{titlesec}
\usepackage{listings}
\usepackage{hyperref}
\usepackage{booktabs}
\usepackage{graphicx}

\definecolor{matd48azul}{HTML}{0B4F6C}
\definecolor{matd48verde}{HTML}{2E8B57}

\titleformat{\section}{\normalfont\Large\bfseries\color{matd48azul}}{\thesection}{1em}{}
\titleformat{\subsection}{\normalfont\large\bfseries\color{matd48azul}}{\thesubsection}{1em}{}

\providecommand{\ListaTitulo}{Lista de Exercícios}

\setlength{\headheight}{14pt}
\pagestyle{fancy}
\fancyhf{}
\lhead{\small MATD48 --- Planejamento de Experimentos A}
\rhead{\small \ListaTitulo}
\cfoot{\thepage}
\renewcommand{\headrulewidth}{0.4pt}

\lstset{
  basicstyle=\ttfamily\small,
  breaklines=true,
  frame=single,
  columns=fullflexible,
  backgroundcolor=\color{gray!8},
  commentstyle=\color{matd48verde},
  keywordstyle=\color{matd48azul}\bfseries,
  language=R,
  extendedchars=true,
  literate=%
    {á}{{\'a}}1 {à}{{\`a}}1 {â}{{\^a}}1 {ã}{{\~a}}1
    {é}{{\'e}}1 {ê}{{\^e}}1 {í}{{\'i}}1 {ó}{{\'o}}1
    {ô}{{\^o}}1 {õ}{{\~o}}1 {ú}{{\'u}}1 {ç}{{\c c}}1
    {Á}{{\'A}}1 {À}{{\`A}}1 {Â}{{\^A}}1 {Ã}{{\~A}}1
    {É}{{\'E}}1 {Ê}{{\^E}}1 {Í}{{\'I}}1 {Ó}{{\'O}}1
    {Ô}{{\^O}}1 {Õ}{{\~O}}1 {Ú}{{\'U}}1 {Ç}{{\c C}}1
}

\hypersetup{colorlinks=true, linkcolor=matd48azul, urlcolor=matd48azul, citecolor=matd48azul}

\newcommand{\cabecalho}[2]{%
  \begin{center}
    {\Large\bfseries\color{matd48azul} #1}\\[0.3em]
    {\large #2}\\[0.3em]
    \rule{\linewidth}{0.6pt}
  \end{center}
  \vspace{0.5em}
}

\newenvironment{questao}[1]{\vspace{1em}\noindent\textbf{Questão #1.}\hspace{0.4em}}{\par}
\newenvironment{solucao}{\vspace{0.3em}\noindent\textit{Solução.}\hspace{0.4em}\color{black!85}}{\par\vspace{0.5em}}


\usepackage[utf8]{inputenc}
\usepackage[T1]{fontenc}
\usepackage[brazil]{babel}
\usepackage{fullpage}
\usepackage{amsmath,amssymb}
\usepackage{enumitem}
\usepackage{xcolor}
\usepackage{fancyhdr}
\usepackage{titlesec}
\usepackage{listings}
\usepackage{hyperref}
\usepackage{booktabs}
\usepackage{graphicx}

\definecolor{matd48azul}{HTML}{0B4F6C}
\definecolor{matd48verde}{HTML}{2E8B57}

\titleformat{\section}{\normalfont\Large\bfseries\color{matd48azul}}{\thesection}{1em}{}
\titleformat{\subsection}{\normalfont\large\bfseries\color{matd48azul}}{\thesubsection}{1em}{}

\providecommand{\ListaTitulo}{Lista de Exercícios}

\setlength{\headheight}{14pt}
\pagestyle{fancy}
\fancyhf{}
\lhead{\small MATD48 --- Planejamento de Experimentos A}
\rhead{\small \ListaTitulo}
\cfoot{\thepage}
\renewcommand{\headrulewidth}{0.4pt}

\lstset{
  basicstyle=\ttfamily\small,
  breaklines=true,
  frame=single,
  columns=fullflexible,
  backgroundcolor=\color{gray!8},
  commentstyle=\color{matd48verde},
  keywordstyle=\color{matd48azul}\bfseries,
  language=R,
  extendedchars=true,
  literate=%
    {á}{{\'a}}1 {à}{{\`a}}1 {â}{{\^a}}1 {ã}{{\~a}}1
    {é}{{\'e}}1 {ê}{{\^e}}1 {í}{{\'i}}1 {ó}{{\'o}}1
    {ô}{{\^o}}1 {õ}{{\~o}}1 {ú}{{\'u}}1 {ç}{{\c c}}1
    {Á}{{\'A}}1 {À}{{\`A}}1 {Â}{{\^A}}1 {Ã}{{\~A}}1
    {É}{{\'E}}1 {Ê}{{\^E}}1 {Í}{{\'I}}1 {Ó}{{\'O}}1
    {Ô}{{\^O}}1 {Õ}{{\~O}}1 {Ú}{{\'U}}1 {Ç}{{\c C}}1
}

\hypersetup{colorlinks=true, linkcolor=matd48azul, urlcolor=matd48azul, citecolor=matd48azul}

\newcommand{\cabecalho}[2]{%
  \begin{center}
    {\Large\bfseries\color{matd48azul} #1}\\[0.3em]
    {\large #2}\\[0.3em]
    \rule{\linewidth}{0.6pt}
  \end{center}
  \vspace{0.5em}
}

\newenvironment{questao}[1]{\vspace{1em}\noindent\textbf{Questão #1.}\hspace{0.4em}}{\par}
\newenvironment{solucao}{\vspace{0.3em}\noindent\textit{Solução.}\hspace{0.4em}\color{black!85}}{\par\vspace{0.5em}}


\begin{document}

\cabecalho{Lista de Exercícios 07}{Contrastes, comparações múltiplas e polinômios ortogonais (Capítulo 3, segunda metade / Aula 07)}

\noindent \textit{Entrega: até a próxima aula. Justifique todas as respostas e mostre os cálculos
intermediários --- nestas listas, a justificativa vale mais que a resposta final. O gabarito
completo está em \texttt{Gabarito07.pdf}.}

\begin{questao}{1} \textbf{(Ciência de dados --- cálculo de um contraste)}

Uma equipe de engenharia web mede o tempo de carregamento (ms) de uma página sob quatro
configurações de scripts de terceiros: \textbf{Nenhum} (controle), \textbf{Analytics},
\textbf{Anúncios} e \textbf{Combo}, com $n=10$ carregamentos independentes em cada uma. As médias
amostrais observadas foram
\[
\bar y_{\text{Nenhum}} = 860, \quad \bar y_{\text{Analytics}} = 940, \quad
\bar y_{\text{Anúncios}} = 1080, \quad \bar y_{\text{Combo}} = 1120 \ \text{(ms)},
\]
com quadrado médio do erro $QME = 4200$ e $gl_{\text{erro}} = 36$. Considere o contraste
$L = 3\mu_{\text{Nenhum}} - \mu_{\text{Analytics}} - \mu_{\text{Anúncios}} - \mu_{\text{Combo}}$,
que compara o grupo controle contra a média dos demais.

\begin{enumerate}[label=(\alph*)]
  \item Escreva o vetor de coeficientes $\boldsymbol\lambda$ deste contraste e verifique que
    $\mathbf 1'\boldsymbol\lambda = 0$.
  \item Calcule $\hat L$.
  \item Calcule $SQ_L$ e a estatística $F = SQ_L/QME$. Compare com o valor crítico
    $F_{0{,}05;\,1,36} \approx 4{,}11$ e conclua.
\end{enumerate}
\end{questao}

\begin{questao}{2} \textbf{(Ciência de dados --- contrastes ortogonais e partição de $SQ_{\text{Trat}}$)}

Usando os mesmos dados da Questão 1, considere o conjunto de três contrastes definidos pela
tabela de coeficientes abaixo.

\begin{center}
\begin{tabular}{lcccc}
\toprule
Contraste & Nenhum & Analytics & Anúncios & Combo \\
\midrule
$L_1$ (Questão 1) & 3 & $-1$ & $-1$ & $-1$ \\
$L_2$ & 0 & 2 & $-1$ & $-1$ \\
$L_3$ & 0 & 0 & 1 & $-1$ \\
\bottomrule
\end{tabular}
\end{center}

\begin{enumerate}[label=(\alph*)]
  \item Verifique que $L_1$, $L_2$ e $L_3$ são mutuamente ortogonais, calculando os três produtos
    internos $\boldsymbol\lambda_1'\boldsymbol\lambda_2$, $\boldsymbol\lambda_1'\boldsymbol\lambda_3$
    e $\boldsymbol\lambda_2'\boldsymbol\lambda_3$.
  \item Calcule $\hat L_2$, $SQ_{L_2}$ e $\hat L_3$, $SQ_{L_3}$ (mesmas fórmulas da Questão 1).
  \item Sabendo que $SQ_{\text{Trat}} = 440\,000$ nesta ANOVA, verifique que
    $SQ_{L_1} + SQ_{L_2} + SQ_{L_3} = SQ_{\text{Trat}}$ (use o $SQ_{L_1}$ calculado na Questão 1).
  \item Interpretando os três testes $F$ individuais, o que você concluiria sobre \emph{onde}
    está a diferença entre as quatro configurações, além de "existe alguma diferença"?
\end{enumerate}
\end{questao}

\begin{questao}{3} \textbf{(Psicologia --- Scheffé vs. teste ingênuo)}

No experimento de dose de cafeína (0, 50, 100, 150, 200 mg; $t=5$; $n=8$ por dose;
$QME = 40$, $gl_{\text{erro}} = 35$), um pesquisador \emph{examina os dados antes} de decidir
qual comparação testar e escolhe, a posteriori, um contraste específico entre grupos de doses
intermediárias. O contraste resultante tem $SQ_L = 185$.

\begin{enumerate}[label=(\alph*)]
  \item Calcule $F = SQ_L/QME$. Testando ao nível $\alpha = 0{,}05$ contra
    $F_{0{,}05;\,1,35} \approx 4{,}12$, o contraste pareceria significativo?
  \item Como o contraste foi escolhido \emph{depois} de examinar os dados (não planejado a
    priori), o procedimento correto é o de Scheffé, que rejeita $H_0$ apenas se
    $F > (t-1)\,F_{0{,}05;\,t-1,\,gl_{\text{erro}}}$. Usando $F_{0{,}05;\,4,35} \approx 2{,}64$,
    calcule o valor crítico de Scheffé e decida se o contraste permanece significativo.
  \item Explique, em suas palavras, por que o teste do item (a) é anticonservador nesta situação
    --- isto é, por que ele rejeita $H_0$ com mais frequência do que o nível nominal de 5\%
    permitiria, se o pesquisador estiver livre para escolher o contraste depois de ver os dados.
\end{enumerate}
\end{questao}

\begin{questao}{4} \textbf{(Dedução --- independência de contrastes ortogonais)}

Considere a parametrização de médias de casela do DCA, $\mathbf Y = \mathbf X\boldsymbol\mu +
\boldsymbol\varepsilon$, balanceada ($n_i = n$ para todo $i$), com
$\boldsymbol\varepsilon \sim N(\mathbf 0, \sigma^2\mathbf I_N)$. Sejam dois contrastes
$L_1 = \boldsymbol\lambda_1'\boldsymbol\mu$ e $L_2 = \boldsymbol\lambda_2'\boldsymbol\mu$, com
estimadores $\hat L_1 = \boldsymbol\lambda_1'\bar{\mathbf y}$ e
$\hat L_2 = \boldsymbol\lambda_2'\bar{\mathbf y}$, em que $\bar{\mathbf y} = (\bar y_{1\cdot},
\ldots, \bar y_{t\cdot})'$.

\begin{enumerate}[label=(\alph*)]
  \item Mostre que $\mathrm{Cov}(\hat L_1, \hat L_2) = \dfrac{\sigma^2}{n}\boldsymbol\lambda_1'
    \boldsymbol\lambda_2$. (Dica: $\bar{\mathbf y}$ tem matriz de covariância
    $\dfrac{\sigma^2}{n}\mathbf I_t$, porque as $t$ médias amostrais são independentes entre si.)
  \item Conclua que, se $\boldsymbol\lambda_1'\boldsymbol\lambda_2 = 0$ (contrastes ortogonais),
    então $\mathrm{Cov}(\hat L_1, \hat L_2) = 0$.
  \item Como $\hat L_1$ e $\hat L_2$ são combinações lineares do vetor conjuntamente normal
    $\bar{\mathbf y}$, o par $(\hat L_1, \hat L_2)$ é conjuntamente normal. Usando o fato de que,
    para vetores conjuntamente normais, covariância zero implica independência estatística,
    conclua que $\hat L_1 \perp \hat L_2$ e, portanto, que $SQ_{L_1}$ e $SQ_{L_2}$ (funções
    determinísticas de $\hat L_1$ e $\hat L_2$, respectivamente) também são independentes.
\end{enumerate}
\end{questao}

\begin{questao}{5} \textbf{(Dedução --- coeficientes do polinômio ortogonal linear, $t=3$)}

Considere um fator quantitativo com $t=3$ níveis igualmente espaçados. Após centralizar e
escalar, os níveis podem ser representados pelos valores codificados $-1, 0, 1$. Procuramos o
contraste linear $\boldsymbol\lambda_L = (c_1, c_2, c_3)$ que:
\begin{enumerate}[label=(\roman*)]
  \item soma zero: $c_1 + c_2 + c_3 = 0$;
  \item é diretamente proporcional ao nível codificado de cada grupo (isto é, $c_i = k \cdot
    (\text{nível codificado}_i)$ para alguma constante $k \neq 0$);
  \item tem coeficientes inteiros de menor módulo possível.
\end{enumerate}

\begin{enumerate}[label=(\alph*)]
  \item Encontre $\boldsymbol\lambda_L$. Mostre que a condição (i) é automaticamente satisfeita
    por qualquer escolha proporcional aos níveis codificados $-1, 0, 1$, e explique por quê.
  \item O contraste quadrático tabulado para $t=3$ é $\boldsymbol\lambda_Q = (1,-2,1)$. Verifique
    que $\boldsymbol\lambda_L$ (encontrado no item a) e $\boldsymbol\lambda_Q$ são ortogonais.
  \item $t=3$ níveis admitem, no máximo, quantos contrastes ortogonais? $\boldsymbol\lambda_L$ e
    $\boldsymbol\lambda_Q$ esgotam essa contagem?
\end{enumerate}
\end{questao}

\begin{questao}{6} \textbf{(Agricultura --- interpretação de polinômios ortogonais)}

Um experimento testa quatro doses de nitrogênio (0, 40, 80, 120 kg/ha) sobre a produtividade do
milho, com $n$ igual em cada dose. A ANOVA decomposta em componentes polinomiais produziu:

\begin{center}
\begin{tabular}{lccccc}
\toprule
Fonte & $gl$ & $SQ$ & $QM$ & $F$ & valor-$p$ \\
\midrule
Dose (total)   & 3 & 812 & --      & --    & --      \\
\ \ Linear     & 1 & 705 & 705{,}0 & 38{,}1 & $<0{,}001$ \\
\ \ Quadrático & 1 & 98  & 98{,}0  & 5{,}3  & 0{,}028 \\
\ \ Cúbico     & 1 & 9   & 9{,}0   & 0{,}5  & 0{,}491 \\
Erro           & 20 & 370 & 18{,}5 & --    & --      \\
\bottomrule
\end{tabular}
\end{center}

\begin{enumerate}[label=(\alph*)]
  \item Verifique que a soma dos três componentes bate com $SQ$ de "Dose (total)".
  \item Quais componentes são estatisticamente significativos ao nível de 5\%? O que isso
    significa, em termos agronômicos, sobre a forma da curva de resposta da produtividade à dose
    de nitrogênio?
  \item Por que faz sentido usar polinômios ortogonais aqui, e não faria sentido usar o mesmo
    procedimento se as quatro doses fossem, por exemplo, quatro \emph{variedades} de milho?
\end{enumerate}
\end{questao}

\end{document}
